
\section{Background}

\subsection{Softmax Attention}

For a single head with $Q, K, V \in \mathbb{R}^{L \times d_h}$:

\begin{equation}
S = \frac{Q K^\top}{\sqrt{d_h}} \in \mathbb{R}^{L \times L}
\end{equation}

\begin{equation}
P = \text{softmax}(S) \in \mathbb{R}^{L \times L}
\end{equation}

\begin{equation}
O = P V \in \mathbb{R}^{L \times d_h}
\end{equation}

\textbf{Complexity.} $S$ and $P$ require $\mathcal{O}(L^2)$ memory. The two 
matmuls require $\mathcal{O}(L^2 d_h)$ time. For $L = 1024$, $d_h = 64$, 
this is $\sim 67$M FLOPs per head per sequence.

\subsection{Sparsity in Attention}

\textbf{Observation.} The softmax distribution $P$ is typically concentrated:

\begin{equation}
\sum_{j \in \text{top-}k} P_{ij} \gg 1 - \epsilon, \quad k \ll L
\end{equation}

\textbf{Empirical evidence.} \textbf{Figure 3.1} shows the cumulative attention mass 
vs ranked keys for a middle GPT-2 layer. Top 10\% of keys capture 
\textbf{85\%} of mass; top 50\% capture \textbf{97\%}.

\textbf{Implication.} Pruning to the top $k = (1 - \rho) L$ keys retains most 
of the output quality, at fraction $1 - \rho$ of the compute.

\subsection{Top-$k$ Selection}

For a score vector $s \in \mathbb{R}^L$, the top-$k$ selection produces:

\begin{equation}
\text{top-}k(s)_j = \begin{cases} s_j & \text{if } s_j \geq \tau \\ -\infty & \text{otherwise} \end{cases}
\end{equation}

where $\tau$ is the $k$-th largest value.

\textbf{Non-differentiability.} The gradient $\partial \text{top-}k / \partial s$ 
is zero almost everywhere. This is problematic for end-to-end training.

\subsection{Soft Top-$k$ Relaxations}

\textbf{Sigmoid mask.} Given $\tau$, we can define:

\begin{equation}
m_j = \sigma\!\left(\frac{s_j - \tau}{T}\right) \in (0, 1)
\end{equation}

so that $m_j \to 1$ for $s_j > \tau$ and $m_j \to 0$ for $s_j < \tau$. 
As $T \to 0$, $m_j \to \mathbb{1}[s_j > \tau]$, recovering hard top-$k$.

\textbf{Application in log space.} Instead of a multiplicative mask, we add a 
log-mask to the scores:

\begin{equation}
\tilde{s}_j = s_j + \log(m_j + \epsilon)
\end{equation}

This preserves the softmax interpretation: $\text{softmax}(\tilde{s}) = 
\text{softmax}(s) \cdot m / \sum m$.

\subsection{Per-Head Sparsity Budgets}

The retention budget $k_h$ for head $h$ is:

\begin{equation}
k_h = \lfloor (1 - \rho_h) L \rfloor
\end{equation}

where $\rho_h \in [0, 1]$ is the head-specific sparsity ratio.

\textbf{Learning $\rho_h$.} We parameterize via a scalar logit:

\begin{equation}
\rho_h = \sigma(\ell_h), \qquad \ell_h \in \mathbb{R}
\end{equation}

This is differentiable w.r.t. $\ell_h$, allowing gradient-based learning of 
the sparsity budget.

\subsection{Chunked Computation}

To bound peak memory, we process the query dimension in chunks of size $C$:

\begin{equation}
O_{[i:i+C]} = \text{softmax}\!\left(\frac{Q_{[i:i+C]} K^\top}{\sqrt{d_h}}\right) V
\end{equation}

\textbf{Peak memory.} $\mathcal{O}(B H C L)$ instead of $\mathcal{O}(B H L^2)$.

For $C = 64$, $B = 8$, $H = 12$, $L = 512$: $64 \times 12 \times 64 \times 512 \times 2 = 50$ MB 
instead of $8 \times 12 \times 512^2 \times 2 = 400$ MB. \textbf{8$$\times$ reduction.}

\subsection{LoRA Recovery}

LoRA (Hu et al., 2021) adds trainable adapters:

\begin{equation}
W_\text{eff} = W_\text{frozen} + \frac{\alpha}{r_L} B A
\end{equation}

with $B = 0$ at initialization. ASA uses LoRA for recovery after sparsity 
injection.

\subsection{Notation Summary}

\begin{lstlisting}

---
\end{lstlisting}